% PROVENANCE (section-level)
%   status      : CODE_FACT unless marked otherwise
%   produced-by : Recon_scripts/2_Reconstruction.py (global tags, ntuple
%                 trees); 5_BBbkg_weights_optuna_minuit.py lines 93, 809-824
%                 (luminosity scale, sample and tree naming); README.md and
%                 docs/ANALYSIS_CONTEXT.md of the analysis repository
%   commit      : ad2f2ec
%   sample      : Proc16 + Prompt16 data, MC16rd; Run 1 + Run 2; e and mu
%   selection   : n/a
%   corrections : n/a
\section{Data, Simulated Samples, and Software Versions}
\label{sec:data}

\subsection{Data samples}
\label{sec:data:data}

The analysis uses Proc16 and Prompt16 data, covering both Run~1 and Run~2, and
is evaluated separately for the $e$ and $\mu$ channels throughout.
The run periods are initially kept separate because
several inputs are run dependent: the PID correction tables
(Sec.~\ref{sec:corrections:recommended}), the fake-$D$ normalization, and the
generic-$B\bar{B}$ corrections themselves (Sec.~\ref{sec:bbbar:rundep}).
For the final $R(D)$ result, we combine the Run 1 and Run 2 datasets.
Both on-resonance ($\Upsilon(4S)$) and off-resonance samples are processed.


\subsection{Monte Carlo (MC) samples}
\label{sec:data:mc}

We use the MC16rd (run-dependent) MC campaign, which was produced with release-08.

\emph{Generic MC}: 
%decays particles according to their measured branching
%fractions, so that it simulates real data as a whole;
The generic MC sample has four
times the data luminosity, so a factor of $0.25$ is applied whenever an
absolute normalization is compared with data.

\emph{Signal MC}: We produce large samples of events (of order $3\times10^{7}$
for a single mode).
Signal samples are generated for this analysis, and have the same
charged-$D$ lepton final state used in the signal reconstruction.
Some of the signal MC samples have $B^{0}\bar{B}^{0}$ events: the $D$ and $D^{*}$ samples are generated
as $B^{0}\bar{B}^{0}$ only, while the $D^{**}$ sample is generated as both
$B^{0}\bar{B}^{0}$ and $B^{+}B^{-}$.  To summarize the signal MC samples:
\begin{itemize}
  \item one $B \to D\tau\nu$ mode (neutral $B$, charged $D$);
  \item one $B \to D^{*}\tau\nu$ mode (neutral $B$, charged $D^{*}$; the
        neutral $D^{*}$ does not decay to a charged $D$ and so does not
        contribute to this final state);
  \item eight $B \to D^{**}\tau\nu$ modes: four with a neutral $B$ and a
        charged $D^{**}$, and four with a charged $B$ and a neutral
        $D^{**}$, since both charge combinations of $D^{**}$ decay to a
        charged $D$.

        \item \textcolor{red}{\bf $B\to D^{**}\ell\nu$ and semileptonic gap modes will be described.}
\end{itemize}
These three signal MC samples are used to obtain the $Pyhf$ fit templates of
Sec.~\ref{sec:fit:templates}. In addition, the $B\to D\tau\nu$ MC sample is used 
as the signal class in the BDT training (Sec.~\ref{sec:mva:classes}).

Generic MC is used for for various purposes including determination of background templates, fit validations,  and checks of control regions.
Since the signal components do not have sufficient statistics, we use dedicated signal MC samples 
for the signal templates in the PYHF fit.



 
\subsection{Software versions}
\label{sec:data:software}

Three release names appear in the analysis chain and they play different
roles, so they are distinguished here explicitly.

Data and simulation are produced and processed with the basf2 \emph{main}
release \texttt{release-08-00-08}.
The main release contains all basf2 packages and is used for data
taking, MC production and data processing.

The analysis-level reconstruction of Sec.~\ref{sec:reconstruction} is run with
the basf2 \emph{light} release \texttt{light-2505-deimos}, a subset of the main
release containing the packages needed for physics analysis.


Release names appearing inside global-tag names --- for example
\texttt{pid\_nn\_release08\_v1} and
\texttt{tracking\_release08\_20260119\_EnergyMomentumScale} --- refer to the
release used when those payloads were prepared.


Global tags used are the analysis global tag returned by
\texttt{getAnalysisGlobaltag()}, \texttt{pid\_nn\_release08\_v1} for the
neural-network identifiers, and
\texttt{tracking\_release08\_20260119\_EnergyMomentumScale} for the track
momentum-scale and energy-loss payloads.


We use the following versions of the offline Python stack --- \texttt{pyhf}: 0.7.6,
\texttt{cabinetry}: 0.6.0, \texttt{iminuit}: 2.26.0, \texttt{LightGBM}: 4.4.0, \texttt{optuna}: 4.5.0,
\texttt{uproot}: 5.6.1 and the Belle II systematics-corrections framework, version 20260922, to produce
the results in this B2note.

\subsection{Ntuple production and offline samples}
\label{sec:data:ntuples}

The analysis is run on mDST files on the GRID. The output ntuples are downloaded onto KEKCC. These ntuples are skimmed using the multi-class BDT. The resulting output is saved into two sets of
ntuples with loose and tight BDT requirements.
The ntuples with loose BDT requirements are used for checking control regions, while the ntuples with tight BDT requirements are used for the PYHF fit and the final analysis.



\subsection{Event weights and normalization}
\label{sec:data:weights}

There are a number of weights used to correct the two-dimensional histograms that are input to the PYHF fit.
There are weights determined by the Performance group to correct the PID efficiencies and fake rates
for charged hadrons and for leptons.
These are corrections for tracking efficiency, momentum scale, charged track energy loss, and slow $\pi^0$ efficiency.

There are a second class of  weights to correct MC physics modeling. 
These are corrections for the
$B\bar{B}$ background and 
continuum background, which are discussed later in the B2NOTE.










